\chapter{Fazit}
Im Versuch der elektrischen Resonanz wurde das Verhalten von elektrischen Schwingkreisen behandelt. Zunächst wurde eine theoretische Beschreibung bereitgestellt. Dadurch können die Phänomene die im Experiment beobachtet wurden verstanden werden, und es kann gleichzeitig eine Überprüfung der theoretischen Vorhersage mit dem Experiment vorgenommen werden. Dazu wurden 3 verschiedene Schwingkreise - ein Reihenschwingkreis, ein Parallelkreis 1. Ordnung und ein Parallelkreis 2. Ordnung - behandelt. Gemessen wurde jeweils die Impedanz der Schwingkreise. Zunächst wurde ein Reihenkreis mit Hilfe eines analogen Oszilloskops vermessen. Die Ergebnisse dieser Messung stimmten sehr gut mit der Theorie überein. Die präzise Anwendbarkeit dieses Verfahren wurde eindrucksvoll demonstriert. Im zweiten Teil des Versuchs wurden die Schwingkreise digital mit Hilfe von Cassy vermessen. Diese Messungen konnten in viel geringerer Zeit durchgeführt werden. Die Ergebnisse der Messung stimmten allerdings leider nicht so gut mit der theoretischen Erwartung überein, wie das beim analogen Oszilloskop der Fall war. Es wurde insbesondere ein scheinbares Offset in der Messung beobachtet.